\section{NumPy}

\begin{frame}[fragile,t,label=cd-numpy]{1. Arrays}
Um \textit{array} do NumPy armazena elementos do mesmo tipo. Diferentemente de uma lista, permite realizar cálculos numéricos diretamente com seus elementos.
\EMASCcodelabel{Exemplo prático}
\begin{lstlisting}
notas = np.array([5.9, 6.7, 7.5, 8.2, 9.1])
print(notas)
print(notas.dtype)   # Tipo dos elementos
print(notas.ndim)    # Numero de dimensoes
print(notas.shape)   # (5,), cinco elementos
print(notas[0])      # Primeiro elemento
print(notas[1:4])    # Elementos de indices 1 a 3
\end{lstlisting}
\begin{EMASCnote}
O índice começa em zero. O índice final de uma fatia não entra no resultado.
\end{EMASCnote}
\end{frame}

\begin{frame}[fragile,t]{1.1. Matrizes}
Um array com duas dimensões pode representar uma matriz. O atributo \texttt{shape} indica o número de linhas e colunas.
\EMASCcodelabel{Horas de estudo e notas de dois estudantes}
\begin{lstlisting}
matriz = np.array([[2.0, 5.9],
                   [3.0, 6.7]])

print(matriz.shape)    # (2, 2)
print(matriz[0, 1])    # 5.9
print(matriz[:, 0])    # Primeira coluna
print(matriz[1, :])    # Segunda linha
\end{lstlisting}
\begin{EMASCnote}
Em uma matriz, \texttt{matriz[i, j]} acessa o elemento da linha \texttt{i} e coluna \texttt{j}.
\end{EMASCnote}
\end{frame}

\begin{frame}[fragile,t]{1.2. Criação de arrays}
\begin{tabularx}{\linewidth}{@{}lX@{}}
\toprule
\textbf{Função} & \textbf{Descrição} \\
\midrule
\texttt{np.zeros(n)} & Array com $n$ zeros. \\
\texttt{np.ones(n)} & Array com $n$ uns. \\
\texttt{np.arange(inicio, fim, passo)} & Passo fixo, sem incluir o fim. \\
\texttt{np.linspace(inicio, fim, n)} & $n$ valores equidistantes, inclui o fim. \\
\bottomrule
\end{tabularx}
\EMASCcodelabel{Exemplos complementares para consulta}
\begin{lstlisting}
print(np.zeros(3))          # [0. 0. 0.]
print(np.ones(3))           # [1. 1. 1.]
print(np.arange(0, 10, 2))  # [0 2 4 6 8]
print(np.linspace(0, 10, 6)) # [0. 2. 4. 6. 8. 10.]
\end{lstlisting}
\end{frame}

\begin{frame}[fragile,t]{1.3. Operações vetorizadas}
Arrays de mesmo formato operam elemento a elemento. Exemplo: horas de estudo dos mesmos três estudantes em duas semanas.
\EMASCcodelabel{Exemplo prático}
\begin{lstlisting}
semana1 = [2, 3, 4]
semana2 = [3, 4, 5]
print(semana1 + semana2)  # [2, 3, 4, 3, 4, 5]
total = np.array(semana1) + np.array(semana2)
print(total)           # [5 7 9], horas de estudo
print(total / 2)       # Media semanal por estudante
print(notas >= 7)      # Comparacao elemento a elemento
\end{lstlisting}
\begin{EMASCnote}
Em listas, \texttt{+} concatena os elementos. Em arrays, soma elemento a elemento. A comparação retorna valores booleanos.
\end{EMASCnote}
\end{frame}

\begin{frame}[fragile,t]{1.4. Medidas estatísticas com NumPy}
As funções do NumPy permitem calcular as medidas estudadas na parte de Estatística.
\EMASCcodelabel{Exemplo prático}
\begin{lstlisting}
print("Média:", np.mean(notas))
print("Mediana:", np.median(notas))
print("Mínimo:", np.min(notas))
print("Máximo:", np.max(notas))
print("Desvio populacional:", np.std(notas, ddof=0))
print("Desvio amostral:", np.std(notas, ddof=1))
\end{lstlisting}
\begin{EMASCnote}
Na variância, \texttt{ddof=0} usa $n$ para descrever a população. Utilize \texttt{ddof=1} (divisor $n-1$) quando a amostra serve para estimar a variância populacional.
\end{EMASCnote}
\end{frame}
